\scanpage{6}\mspage{\handwritten{5}}

éventuellement muni de la sous-catégorie pleine de tous les systèmes \unsure{locaux} sur $X$, et du foncteur cohomologique sur $D^b_{lc}(X)$ à valeurs dqns \unsure{laddtë|laddte} catégorie, et bien sûr du produit tensoriel (mais alors on sort de $D^b$ pour entrer dans $D^-$, redactor demerdetur).

II \quad $n$-catégories, catégories $n$-uples, et $(\mathrm{Gr})$-catégories.

Appelons bi-catégorie une catégorie dans$(\mathrm{Cat})$, i.e. un triple $(C^{\deleted{o}}_{\handwritten{0}},C^{\deleted{1}}_{\handwritten{1}},\ \deleted{\unsure{)\text{ où }}}p_1,p_2,\handwritten{\pi})$, où $C^o$ est une catégorie, $p_1,p_2:C^{\deleted{1}}_{\handwritten{1}}\handwritten{\rightrightarrows}C^{\deleted{o}}_{\handwritten{0}}$ deux foncteurs, $\handwritten{\pi}:(C^{\deleted{1}}_{\handwritten{1}},p_1)\mathbin{x}_{C^o}(C^{\deleted{1}}_{\handwritten{1}},p_2)\handwritten{\to}C^{\deleted{1}}_{\handwritten{1}}$ un foncteur, avec les axiomes habituels. Interprétant $(\mathrm{Cat})$ comme la sous-catégorie pleine de $(\mathrm{Ssimpl})$ des objets qui .... , et de même pour les catégories à valeurs dans une catégorie telle $(\mathrm{Cat})$, on trouve que la catégorie \inserted{$(\mathrm{Bicat})$} des $(\mathrm{Cat})$-catégories est équivalente à la sous-catégorie pleine de la catégorie des ensembles bi-semisimpliciaux

\diagram{\(\begin{array}{cccc}
X^{\deleted{oo}}_{\handwritten{00}} & X^{\deleted{o1}}_{\handwritten{01}} & \ldots & X^{\deleted{o,n}}_{\handwritten{0n}}\\
X^{\deleted{1o}}_{\handwritten{10}} & X^{\deleted{11}}_{\handwritten{11}} & \ldots & X^{\deleted{1,n}}_{\handwritten{1n}}\\
\multicolumn{2}{l}{.......} & &
\end{array}\)}

formée des objets qui ligne par ligne et colonne par colonne \unsure{satisfoht|satisfont} la condition d'exctitude familière (transformer sommes amalgamées en produits fibrés). De même la catégorie \deleted{des} $(\mathrm{Tricat})$ des $(\mathrm{Bicat})$-catégories \unsure{ssinterprète|s'interprète} en termes d'objets tri-\inserted{semi}simpliciaux, etc.

Exemple: soit $C$ une catégorie, posons $C^{\deleted{o}}_{\handwritten{0}}=C$, $C^{\deleted{1}}_{\handwritten{1}}=\underline{\mathrm{Fl}}(C)$, avec structures $p_1,p_2$, évidentes (d'où une $(\mathrm{Cat})$-catétorie dont les composantes semi-simpliciales supérieures sont les catégories $\underline{\mathrm{Fl}}^n(C)$ des $n$-flêches de $C$ ...). L'ensemble bi-semisimplicial correspondant a pour composantes les \unsure{ensembles} $\mathrm{Hom}(\handwritten{\Delta}^n\mathbin{x}\handwritten{\Delta}^m,C)$. Par exemple, les objets de $X^{\deleted{11}}_{\handwritten{11}}$ sont les carrés commutatifs de flêches de $C$, qu'on peut interpréter (de deux manières différentes) \unsure{comme}

% archive sheet number: 5
